% ======================================================================
%  Memoria del TFG: Automatización del Procesamiento de Boletines
%  Oficiales y Consulta en Lenguaje Natural para Entidades Locales
%  Mediante LLM. Plantilla UMA ETSI Informática (ver README.txt).
%
%  Estructura:
%    preliminares/  resumen, abstract, agradecimientos, acrónimos
%    capitulos/     un fichero por capítulo
%    apendices/     un fichero por apéndice
%    figuras/       imágenes y diagramas
%  Compilar con `latexmk` (ver .latexmkrc: XeLaTeX + biber + nomencl).
%  En Overleaf: compilador XeLaTeX.
% ======================================================================
\documentclass[12pt,a4paper,twoside,openright]{book}

%───────── Auxiliary packages ──────────────────────────────────────────
\usepackage{csquotes}
\usepackage{graphicx}
\graphicspath{{figuras/}}
\usepackage{booktabs}
\usepackage{tabularx}
\usepackage{siunitx}
\usepackage{amsmath,amssymb}
\usepackage{algorithm}
\usepackage{algpseudocode}
\usepackage{listings}
\usepackage{nomencl}
\makenomenclature
\usepackage[
  backend=biber,
  style=numeric,
  sorting=nyt
]{biblatex}
\addbibresource{references.bib}

% Listings config
\lstset{
  basicstyle=\ttfamily\small,
  numbers=left,
  numberstyle=\tiny,
  columns=fullflexible,
  frame=single,
  breaklines=true,
  tabsize=2
}

%───────── Template, covers, back cover ────────────────────────────────
\usepackage{template/sty_files/template-TFGspanish-uma}

\usepackage{template/sty_files/cotutor_cover-blueuma}
\usepackage{template/sty_files/cotutor_cover-whiteuma}

\usepackage{template/sty_files/backcover-uma}

% Datos comunes a las dos portadas
\newcommand{\tfgTitleES}{Automatización del Procesamiento de Boletines Oficiales y Consulta en Lenguaje Natural para Entidades Locales Mediante LLM}
\newcommand{\tfgTitleEN}{Automation of Official Gazette Processing and Natural Language Query Services for Local Authorities Using LLM}
\newcommand{\tfgAuthor}{Antonio Bermúdez Rodríguez}
\newcommand{\tfgTutor}{Dr.~Cristóbal Barba González}
\newcommand{\tfgCotutor}{Sandro José Hurtado Requena}  % TODO: añadir "Dr." si procede
\newcommand{\tfgDept}{Departamento de Lenguajes y Ciencias de la Computación}  % TODO: confirmar departamento
\newcommand{\tfgCityDate}{Málaga, Diciembre, 2026}  % TODO: fecha de entrega

%======================================================================
\begin{document}

%────────────────────── Blue front cover (Cotutor) ───────────────────────
\MakeBlueUMACover{
  degree     = {Grado en Ingeniería Informática},
  titleES    = {\tfgTitleES},
  titleEN    = {\tfgTitleEN},
  author     = {\tfgAuthor},
  tutor      = {\tfgTutor},
  cotutor    = {\tfgCotutor},
  dept       = {\tfgDept},
  cityDate   = {\tfgCityDate},
  logoLeft   = {template/logos/NEG-uma-logo.png},
  logoRight  = {template/logos/NEG-etsi-logo.png}
}

%────────────────────── White front cover (Cotutor) ───────────────────────
\MakeWhiteUMACover{
  degree     = {Graduado en Ingeniería Informática},
  titleES    = {\tfgTitleES},
  titleEN    = {\tfgTitleEN},
  author     = {\tfgAuthor},
  tutor      = {\tfgTutor},
  cotutor    = {\tfgCotutor},
  dept       = {\tfgDept},
  cityDate   = {\tfgCityDate},
  defenseDate = {},  % TODO: fecha de defensa
  logoLeft   = {template/logos/POS-uma-logo.png},
  logoRight  = {template/logos/POS-etsi-logo.png}
}

\frontmatter

%────────────────────── Resumen, abstract, agradecimientos ─────────────
% Resumen (150--300 palabras) y palabras clave (3--5), en castellano.
% TODO: problema, solución, estudio de chunking y resultado principal.
\renewcommand{\abstractname}{Resumen}
\begin{abstract}
% TODO: redactar el resumen.

\medskip
\noindent\textbf{Palabras clave:} RAG, LLM, boletines oficiales, administración local, \emph{chunking}
\end{abstract}

% Abstract (150--300 words) and keywords (3--5), in English.
\renewcommand{\abstractname}{Abstract}
\begin{abstract}
% TODO: translate the Spanish summary.

\medskip
\noindent\textbf{Keywords:} RAG, LLM, official gazettes, local government, chunking
\end{abstract}

% Agradecimientos (opcional).
\begin{umaacknowledgments}
% TODO: agradecimientos.
\end{umaacknowledgments}


%────────────────────── TOC and lists ──────────────────────────────────
\clearpage
\tableofcontents
\clearpage
\listoffigures
\clearpage
\listoftables
\clearpage

%────────────────────── Acrónimos ──────────────────────────────────────
\renewcommand{\nomname}{Acrónimos}
% Glosario de acrónimos (paquete nomencl). Añadir aquí cada acrónimo nuevo.
\nomenclature{ANN}{\emph{Approximate Nearest Neighbor}}
\nomenclature{BOE}{Boletín Oficial del Estado}
\nomenclature{BOJA}{Boletín Oficial de la Junta de Andalucía}
\nomenclature{DLQ}{\emph{Dead Letter Queue}}
\nomenclature{HNSW}{\emph{Hierarchical Navigable Small World}}
\nomenclature{LLM}{\emph{Large Language Model}}
\nomenclature{OTel}{OpenTelemetry}
\nomenclature{RAG}{\emph{Retrieval-Augmented Generation}}
\nomenclature{RRF}{\emph{Reciprocal Rank Fusion}}
\nomenclature{SDD}{\emph{Spec-Driven Development}}
\nomenclature{TTFT}{\emph{Time To First Token}}
\printnomenclature

\clearpage

\mainmatter
%────────────────────── Capítulos ──────────────────────────────────────
% ======================================================================
% Capítulo: Introducción
% Contexto, motivación, objetivos (OE1--OE7), alcance y estructura. Extensión orientativa: 4--5 páginas.
% ======================================================================
\chapter{Introducción}
\label{cap:introduccion}

\section{Contexto y problema}
\label{sec:contexto}

El ordenamiento jurídico español obliga a las administraciones públicas a mantenerse al día de un flujo constante de normas, subvenciones, decretos y resoluciones. Existe, sin embargo, una brecha entre la velocidad a la que se produce esa normativa y la capacidad de las entidades locales para absorberla y entenderla.

En un ayuntamiento, este seguimiento se hace hoy de forma manual, fragmentada y descentralizada: el personal técnico revisa cada día varios diarios oficiales, extensos y redactados en un lenguaje jurídico complejo. Esta carga administrativa supone una ineficiencia operativa notable y, sobre todo, eleva el riesgo de error humano y de fatiga informativa. El problema es más agudo en los municipios pequeños y medianos, que, a diferencia de los ministerios o de las capitales autonómicas, trabajan con plantillas y presupuestos limitados.

Este trabajo aborda ese problema en el ámbito de la Comunidad Autónoma de Andalucía, con dos fuentes: el Boletín Oficial del Estado (BOE) y el Boletín Oficial de la Junta de Andalucía (BOJA). La propuesta es una plataforma que lee esos boletines de forma automática cada día y permite al personal municipal consultarlos en lenguaje natural, como una herramienta de apoyo a la decisión y de mejora de sus flujos de trabajo.

Para ello se recurre a los modelos de lenguaje de gran escala (LLM), que, apoyados en la arquitectura Transformer \cite{vaswani2017}, ofrecen una capacidad sin precedentes para procesar y comprender textos extensos con estructuras semánticas complejas. Pero el dominio exige un rigor legal estricto, y estos modelos pueden generar respuestas plausibles pero falsas, las llamadas \emph{alucinaciones} \cite{huang2023}. Por eso el sistema no responde de memoria: recupera primero los fragmentos de los boletines relevantes para cada pregunta y genera la respuesta a partir de ellos, citándolos (\emph{Retrieval-Augmented Generation}, RAG \cite{lewis2020}). Esta solución se apoya en una base de datos vectorial para la recuperación eficiente de información \cite{pan2024survey} y en una herramienta de orquestación de IA \cite{semantickernel}.
% NOTA: el anteproyecto citaba aquí Wang et al. (2024), referencia que no se ha podido localizar; se sustituye por el survey de Pan, Wang y Li.

\section{Motivación}
\label{sec:motivacion}

La primera motivación es práctica. Una convocatoria de subvenciones que pasa desapercibida o un plazo legal que se incumple sin querer tienen un coste real para un ayuntamiento y para sus vecinos. Automatizar la lectura diaria de los boletines y permitir preguntar por ellos en lenguaje natural reduce ese riesgo y libera tiempo del personal técnico para tareas de más valor.

La segunda motivación es de ingeniería del software. Los sistemas que integran modelos de IA plantean retos que el software tradicional no tiene \cite{ozkaya2020}: su comportamiento no es determinista, dependen de proveedores externos con cuotas y costes variables, y su calidad no se puede comprobar solo con pruebas unitarias. Este trabajo ha sido una oportunidad para abordar esos retos con una arquitectura distribuida y una metodología guiada por especificaciones (capítulo~\ref{cap:metodologia}).

A medida que avanzó el desarrollo, y en especial tras mostrar una primera versión a personal municipal, aparecieron necesidades que el planteamiento inicial no recogía y que han dado forma a parte del sistema:

\begin{itemize}
  \item \textbf{Ninguna espera innecesaria y ninguna pérdida de datos.} La ingesta de un boletín implica descargarlo, fragmentarlo, generar sus \emph{embeddings} y, en la pipeline jerárquica, llamar a un modelo de lenguaje por cada fragmento. Hacerlo de forma síncrona bloquearía a quien lo lanza y perdería el trabajo ante cualquier fallo. Por eso todo el procesamiento es asíncrono y pasa por colas de mensajes persistentes, con reintentos y una cola de mensajes fallidos: el usuario no espera a que termine y ningún documento se pierde por un error transitorio. El mismo principio se aplica al chat, donde los mensajes se guardan en segundo plano para que la respuesta no espere a la base de datos.
  \item \textbf{Una búsqueda rápida con modelos de \emph{embeddings} mayores.} Los índices aproximados de pgvector no admiten vectores de más de 2.000 dimensiones, lo que impedía usar con índice modelos de \emph{embeddings} de mayor dimensión. % VERIFICAR: límite de 2.000 dimensiones de los índices HNSW e IVFFlat de pgvector (documentado en docs/adr/0001-vector-index-dimensions.md; contrastar con la documentación de pgvector).
  Los vectores se trasladaron a Qdrant, una base de datos vectorial dedicada, y PostgreSQL se mantuvo como fuente de verdad para el resto de datos y para la búsqueda léxica.
  \item \textbf{Un asistente que conoce a quien pregunta.} El historial de conversaciones permite retomar consultas anteriores, y su análisis sugiere a cada usuario los temas y municipios que le interesan, que después se usan para adaptar las respuestas.
  \item \textbf{Elegir con datos qué estrategia funciona mejor.} La Arena compara a ciegas las respuestas de dos pipelines y deja que el personal municipal vote cuál prefiere, y las métricas de coste muestran cuánto cuesta cada respuesta según el modelo y la técnica empleados.
  \item \textbf{Varios ayuntamientos en la misma plataforma.} Cada ayuntamiento es un \emph{tenant} aislado, con sus usuarios, conversaciones y configuración de IA propios, de modo que la plataforma puede abrirse a nuevas entidades sin trabajo adicional. Además, está preparada para que cada ayuntamiento use su propia clave de API en el modelo que redacta las respuestas del chat, mientras la plataforma se encarga de la ingesta y de los \emph{embeddings}.
\end{itemize}

% TODO (Antonio): ampliar cada punto con el contexto concreto en que surgió (por ejemplo, el feedback recibido sobre la primera versión).

\section{Objetivos}
\label{sec:objetivos}

\subsection*{Objetivo general}

El objetivo general de este trabajo es diseñar, implementar y desplegar una plataforma conversacional que automatice la lectura diaria del Boletín Oficial del Estado (BOE) y del Boletín Oficial de la Junta de Andalucía (BOJA) y que permita al personal de las entidades locales consultar su contenido en lenguaje natural. Las respuestas se generan mediante \emph{Retrieval-Augmented Generation} (RAG) \cite{lewis2020}, de forma que cada respuesta se apoya en fragmentos concretos de los boletines y los cita, lo que reduce el riesgo de alucinaciones \cite{huang2023}.

\subsection*{Objetivos específicos}

El objetivo general se descompone en siete objetivos específicos, que se corresponden con los entregables del anteproyecto. Cada objetivo se ha concretado después en una o varias \emph{capacidades} de OpenSpec, la herramienta de especificación que se describe en la sección~\ref{sec:sdd}. Las especificaciones de esas capacidades, que se encuentran en el directorio \texttt{openspec/specs} del repositorio, son la fuente de las funcionalidades y requisitos del capítulo~\ref{cap:requisitos}.

{\sloppy
\begin{description}
  \item[OE1. Ingesta automática de boletines.] Desarrollar un microservicio en Go que descargue cada día los documentos del BOE y del BOJA, aplique reglas de filtrado configurables, guarde los originales y el texto estructurado en un almacenamiento compatible con S3 y encole cada documento para su procesamiento. Capacidades: \texttt{go-scraper}, \texttt{scraper-dynamic-filtering} y \texttt{messaging-reliability}.

  \item[OE2. Orquestación del RAG en .NET.] Construir un \emph{backend} en .NET que fragmente los documentos, genere sus \emph{embeddings}, recupere los fragmentos relevantes para cada pregunta y genere la respuesta con Semantic Kernel \cite{semantickernel}, con un proveedor de modelos de lenguaje intercambiable. Capacidades: \texttt{chunking-embeddings}, \texttt{hierarchical-retrieval}, \texttt{query-api-semantic-kernel}, \texttt{user-api-keys-vault} y \texttt{llm-provider-security}.

  \item[OE3. Interfaz conversacional en tiempo real.] Ofrecer un \emph{frontend} en Angular con un chat que muestre la respuesta a medida que se genera, con las citas a los boletines de origen, el historial de conversaciones y una navegación adaptada al móvil. Capacidades: \texttt{signalr-integration}, \texttt{chat-ui-connection}, \texttt{chat-history}, \texttt{mobile-navigation} y \texttt{session-timeout}.

  \item[OE4. Modelo de datos híbrido.] Diseñar un modelo de datos que combine una base de datos relacional (PostgreSQL) como fuente de verdad con una base de datos vectorial (Qdrant) con indexación aproximada de vecinos más cercanos (ANN), y que mantenga ambas coherentes. Capacidades: \texttt{qdrant-vector-store} e \texttt{ingestion-consistency}.

  \item[OE5. Estudio comparativo de estrategias de fragmentación.] Comparar varias estrategias de fragmentación (\emph{chunking}) de los boletines y medir su efecto en la calidad de las respuestas y en el coste de ingesta. Capacidades: \texttt{bulk-reprocessing-pipeline}, \texttt{ingestion-telemetry} y \texttt{question-arena}. El diseño y los resultados del estudio se presentan en el capítulo~\ref{cap:estudio-chunking}.

  \item[OE6. Orquestación y despliegue.] Orquestar todos los servicios con .NET Aspire durante el desarrollo y desplegarlos con Docker Compose en producción, con integración y despliegue continuos y observabilidad de extremo a extremo. Capacidades: \texttt{build-quality-gates}, \texttt{observability-grafana} y \texttt{clean-architecture-boundaries}.

  \item[OE7. Validación con casos reales.] Preparar la plataforma para su uso por varios ayuntamientos a la vez, con los datos de cada uno aislados, y disponer de los instrumentos necesarios para validarla con personal municipal. Capacidades: \texttt{user-auth-auth0}, \texttt{multi-tenant-user-profiles}, \texttt{api-authorization}, \texttt{question-arena} y \texttt{admin-metrics-dashboard}.
\end{description}
\par}

El cumplimiento de cada objetivo se revisa en el capítulo~\ref{cap:conclusiones}.

\section{Alcance y limitaciones}
\label{sec:alcance}

\paragraph{Fuentes.} El sistema cubre dos boletines: el BOE, de ámbito estatal, y el BOJA, de ámbito autonómico andaluz. Se descargan cada día de forma automática y los administradores de la plataforma pueden decidir, con reglas de filtrado por departamento, sección, palabra clave o canal, qué documentos se procesan. Quedan fuera los boletines oficiales de las provincias, los de otras comunidades autónomas y la normativa propia de cada ayuntamiento (ordenanzas y acuerdos de pleno), Los boletines provinciales se descartaron por la dificultad de extraer sus datos de forma automática: muchos protegen el acceso con CAPTCHA y publican sus anuncios solo en PDF, sin una API ni un formato estructurado como el XML del BOE. La arquitectura de ingesta permite, de todos modos, añadir nuevas fuentes implementando un nuevo procesador sin cambiar el resto del sistema, por lo que su incorporación queda como línea futura.

\paragraph{Usuarios.} Los destinatarios son el personal de las entidades locales: funcionarios y personal administrativo que necesitan localizar y entender las normas, convocatorias y plazos que afectan a su municipio. La plataforma es \emph{multi-tenant}: cada ayuntamiento tiene sus usuarios, conversaciones y configuración aislados, y hay perfiles de administrador de la organización y de la plataforma. Aunque las respuestas se redactan en lenguaje claro, el sistema no está pensado como un servicio abierto a la ciudadanía.

\paragraph{Funcionalidad.} El trabajo abarca el ciclo completo: la ingesta de los boletines, su fragmentación e indexación, la recuperación y generación de respuestas con RAG, la interfaz de chat en tiempo real, la comparación de \emph{pipelines} con la Arena, el control de coste por consulta y el despliegue con observabilidad. Incluye además el estudio comparativo de estrategias de fragmentación previsto en el anteproyecto (capítulo~\ref{cap:estudio-chunking}).

\paragraph{Mitigación de alucinaciones.} La principal salvaguarda frente a las alucinaciones es que el modelo responde únicamente a partir de los fragmentos recuperados, cita en cada respuesta los documentos en que se basa con su referencia y su enlace al boletín, y tiene instrucciones de reconocer que no dispone de la información cuando los documentos no la contienen (sección~\ref{sec:alucinaciones}).

\paragraph{Limitaciones.} Conviene dejar claros desde el principio los límites del sistema:
\begin{itemize}
  \item \textbf{No sustituye el criterio profesional.} RAG reduce las alucinaciones pero no las elimina: incluso herramientas jurídicas comerciales basadas en recuperación siguen cometiendo errores en una parte apreciable de las consultas \cite{magesh2025}. Las respuestas son un apoyo para localizar y entender la información, no asesoramiento jurídico, y el usuario debe comprobar la fuente citada antes de tomar una decisión.
  \item \textbf{La calidad depende de lo indexado.} El sistema solo conoce los documentos que se han descargado y procesado; una norma que no se ha ingerido, o que las reglas de filtrado han descartado, no aparecerá en las respuestas.
  \item \textbf{Dependencia de proveedores de modelos.} En producción, la generación de respuestas depende de proveedores externos, con cuotas, latencias y costes que no controla el proyecto; por eso el sistema permite cambiar de proveedor por configuración y registra el coste de cada respuesta. En local se usan modelos abiertos con Ollama, cuya calidad es menor.
  \item \textbf{Validación limitada.} Una primera versión se mostró a personal municipal y su opinión orientó parte del desarrollo, pero la validación formal con funcionarios en uso real depende de que las entidades concedan el acceso. La plataforma incluye los instrumentos para hacerla (la Arena, los perfiles por organización y las métricas), y su alcance real se describe en la sección~\ref{sec:validacion-usuarios}.
  \item \textbf{Idioma.} El sistema está pensado para documentos y preguntas en castellano; no se han tratado las lenguas cooficiales ni otros idiomas.
\end{itemize}

\section{Estructura de la memoria}
\label{sec:estructura}

El resto de la memoria se organiza en los siguientes capítulos:

\begin{itemize}
  \item El \textbf{capítulo~\ref{cap:estado-arte}} repasa los fundamentos en los que se apoya el trabajo: los modelos de lenguaje y la arquitectura Transformer, las alucinaciones, la generación aumentada por recuperación, las estrategias de fragmentación, las bases de datos vectoriales, la recuperación híbrida, la evaluación de sistemas RAG y los marcos de orquestación. Termina con los trabajos relacionados y la justificación de las tecnologías empleadas.
  \item El \textbf{capítulo~\ref{cap:metodologia}} describe la metodología: Scrum adaptado a un desarrollo individual, el desarrollo guiado por especificaciones con OpenSpec, la gestión del código y la calidad, el uso de la inteligencia artificial durante el desarrollo y la planificación temporal.
  \item El \textbf{capítulo~\ref{cap:requisitos}} recoge el análisis de requisitos: los actores, los requisitos funcionales y no funcionales y los casos de uso.
  \item El \textbf{capítulo~\ref{cap:diseno}} presenta el diseño del sistema: la arquitectura general, el modelo de datos híbrido, la indexación vectorial, el flujo de ingesta asíncrona, las \emph{pipelines} RAG y la abstracción del proveedor de IA.
  \item El \textbf{capítulo~\ref{cap:implementacion}} detalla la implementación de cada componente: el módulo de ingesta en Go, el \emph{backend} .NET, la orquestación RAG, la Arena, el \emph{frontend} Angular, el despliegue y la observabilidad.
  \item El \textbf{capítulo~\ref{cap:estudio-chunking}} contiene el estudio comparativo de estrategias de fragmentación: las estrategias, el diseño experimental, las métricas y los resultados.
  \item El \textbf{capítulo~\ref{cap:pruebas}} describe las pruebas y la validación del sistema, incluida la validación con usuarios reales y el rendimiento.
  \item El \textbf{capítulo~\ref{cap:conclusiones}} revisa el cumplimiento de los objetivos, resume las aportaciones y las lecciones aprendidas y propone líneas futuras.
\end{itemize}

Los apéndices recogen los resultados completos del estudio de fragmentación (apéndice~\ref{ap:chunking}), el manual de usuario (apéndice~\ref{ap:manual-usuario}), la guía de despliegue (apéndice~\ref{ap:despliegue}), el contrato de mensajería y las capacidades de OpenSpec (apéndice~\ref{ap:especificaciones}) y los \emph{prompts} empleados (apéndice~\ref{ap:prompts}).

% ======================================================================
% Capítulo: Estado del arte y tecnologías
% Fundamentos de LLM y RAG, chunking, bases vectoriales, evaluación y tecnologías. Extensión orientativa: 7--9 páginas.
% ======================================================================
\chapter{Estado del arte y tecnologías}
\label{cap:estado-arte}

\section{Modelos de lenguaje y arquitectura Transformer}
\label{sec:llm}
% TODO: Capacidades, ventana de contexto y coste \cite{vaswani2017}.

\section{Alucinaciones en LLM}
\label{sec:alucinaciones}
% TODO: Taxonomía y por qué son críticas en el dominio jurídico \cite{huang2023}.

\section{Retrieval-Augmented Generation}
\label{sec:rag}
% TODO: Esquema general y variantes: \emph{naive}, \emph{advanced}, jerárquico \emph{small-to-big} \cite{lewis2020,gao2023}.

\section{Estrategias de fragmentación (\emph{chunking})}
\label{sec:chunking}
% TODO: Fixed-size con solapamiento, recursive character splitting, semantic chunking y padre/hijo.

\section{Embeddings y bases de datos vectoriales}
\label{sec:vectoriales}
% TODO: Similitud coseno, ANN (HNSW, IVFFlat), pgvector frente a Qdrant \cite{wang2024,malkov2020}.

\section{Recuperación híbrida}
\label{sec:hibrida}
% TODO: Búsqueda densa, búsqueda léxica (\texttt{tsvector}) y fusión RRF \cite{cormack2009}.

\section{Evaluación de sistemas RAG}
\label{sec:evaluacion-rag}
% TODO: RAGAS \cite{es2023ragas}, TruLens, métricas de fidelidad y contexto; evaluación por preferencia humana \cite{chiang2024}.

\section{Orquestación de IA}
\label{sec:orquestacion}
% TODO: Semantic Kernel \cite{semantickernel} frente a LangChain y LlamaIndex.

\section{Trabajos relacionados}
\label{sec:relacionados}
% TODO: Asistentes de IA en la administración pública y herramientas jurídicas; qué aporta este TFG.

\section{Tecnologías empleadas}
\label{sec:tecnologias}
% TODO: Tabla con cada tecnología y el motivo de su elección (ver tabla de ejemplo abajo).

\begin{table}[ht]
  \centering
  \caption{Tecnologías empleadas y motivo de su elección.}
  \label{tab:tecnologias}
  \begin{tabularx}{\linewidth}{lX}
    \toprule
    Tecnología & Motivo de la elección \\
    \midrule
    .NET 10 y Aspire & % TODO
    \\
    Go & % TODO
    \\
    Semantic Kernel & % TODO
    \\
    PostgreSQL y Qdrant & % TODO
    \\
    RabbitMQ & % TODO
    \\
    Angular & % TODO
    \\
    \bottomrule
  \end{tabularx}
\end{table}

% ======================================================================
% Capítulo: Metodología y planificación
% Scrum individual + SDD con OpenSpec, calidad y planificación real frente a la del anteproyecto. Extensión orientativa: 3--4 páginas.
% ======================================================================
\chapter{Metodología y planificación}
\label{cap:metodologia}

\section{Scrum adaptado a un desarrollo individual}
\label{sec:scrum}
% TODO: Sprints de 2--3 semanas, backlog y revisión \cite{schwaber2020}.

\section{Desarrollo guiado por especificaciones}
\label{sec:sdd}
% TODO: SDD para IA \cite{ramos2026} concretado con OpenSpec: propuesta, diseño, tareas y archivo. Ejemplos de specs.

\section{Gestión del código y calidad}
\label{sec:calidad}
% TODO: Git, PR por cambio, CI obligatoria, agentes de IA con revisión humana.

\section{Planificación temporal}
\label{sec:planificacion}
% TODO: Las 9 fases y 296 h del anteproyecto frente a la dedicación real (Gantt jul.--oct. 2026) y desviaciones.

\section{Recursos}
\label{sec:recursos}
% TODO: VS Code, Ollama, Podman, servicios en la nube y coste de inferencia.

% ======================================================================
% Capítulo: Análisis de requisitos
% Fase 1 del anteproyecto. Extensión orientativa: 4--5 páginas.
% ======================================================================
\chapter{Análisis de requisitos}
\label{cap:requisitos}

\section{Actores}
\label{sec:actores}

Los actores se corresponden con los roles definidos en Auth0, donde cada ayuntamiento es una \emph{organización} (un \emph{tenant} de la plataforma) y los roles viajan en el token de acceso:

\begin{itemize}
  \item \textbf{Usuario municipal.} Es el usuario principal: un funcionario o un miembro del personal administrativo de un ayuntamiento que consulta los boletines a través del chat, revisa su historial, mantiene sus temas de interés y vota en la Arena. Solo ve los datos de su propio ayuntamiento.
  \item \textbf{Administrador de organización} (rol \texttt{admin-organizacion}). Además de lo anterior, configura el modelo de IA de su ayuntamiento y su clave de API, consulta el consumo de IA de su organización y accede a la página de la Arena.
  \item \textbf{Administrador de plataforma} (rol \texttt{admin-plataforma}). Gestiona la base de conocimiento compartida: las reglas de filtrado del \emph{scraper}, las descargas bajo demanda, el reprocesado de documentos y el coste de ingesta de toda la plataforma.
  \item \textbf{Sistema programado.} El planificador del \emph{scraper}, que lanza cada día la descarga de los boletines del día anterior sin intervención humana.
\end{itemize}

Los documentos de los boletines y sus vectores son globales, de modo que todos los ayuntamientos consultan la misma base de conocimiento. En cambio, las conversaciones, los votos de la Arena, las preferencias y la configuración de IA pertenecen a un ayuntamiento y a un usuario concretos.

\section{Requisitos funcionales}
\label{sec:rf}

Los requisitos funcionales se han extraído de las especificaciones de OpenSpec (\texttt{openspec/specs}), que describen el comportamiento del sistema tal como está implementado. Cada requisito indica la capacidad de la que procede, de modo que puede consultarse allí con sus escenarios de aceptación. Se agrupan en cinco áreas, recogidas en las tablas~\ref{tab:rf-ingesta} a~\ref{tab:rf-admin}.

\begin{table}[htbp]
\centering\small
\caption{Requisitos funcionales de ingesta de boletines}
\label{tab:rf-ingesta}
\begin{tabularx}{\textwidth}{@{}l X l@{}}
\toprule
\textbf{Id.} & \textbf{Requisito} & \textbf{Capacidad} \\
\midrule
RF1 & Descargar cada día los documentos del BOE y del BOJA del día anterior de forma automática. & \texttt{go-scraper} \\
RF2 & Lanzar descargas bajo demanda para un rango de fechas (31 días como máximo), sin permitir dos trabajos a la vez. & \texttt{go-scraper} \\
RF3 & Descartar los documentos que no cumplen las reglas de filtrado, que se guardan en la base de datos y se editan desde la interfaz. Si las reglas no se pueden obtener, no se descarga nada. & \texttt{scraper-dynamic-filtering} \\
RF4 & Guardar el original y el texto estructurado de cada documento en el almacenamiento S3 antes de encolarlo para su procesamiento. & \texttt{go-scraper} \\
RF5 & Fragmentar cada documento, generar los \emph{embeddings} de sus fragmentos y guardarlos en PostgreSQL y Qdrant. & \texttt{chunking-embeddings} \\
RF6 & Procesar cada documento con dos pipelines: la \emph{baseline} (fragmentos de tamaño fijo con solapamiento) y la jerárquica (documentos padre e hijos enriquecidos con contexto). & \texttt{hierarchical-retrieval} \\
RF7 & Reprocesar en bloque los documentos almacenados por una o por las dos pipelines, sin duplicar fragmentos ni dejar vectores huérfanos. & \texttt{bulk-reprocessing-pipeline} \\
\bottomrule
\end{tabularx}
\end{table}
\begin{table}[htbp]
\centering\small
\caption{Requisitos funcionales de consulta}
\label{tab:rf-consulta}
\begin{tabularx}{\textwidth}{@{}l X l@{}}
\toprule
\textbf{Id.} & \textbf{Requisito} & \textbf{Capacidad} \\
\midrule
RF8 & Responder a preguntas en lenguaje natural a partir de los fragmentos recuperados de los boletines. & \texttt{query-api-semantic-kernel} \\
RF9 & Mostrar la respuesta a medida que se genera. & \texttt{signalr-integration} \\
RF10 & Citar en cada respuesta los documentos usados, con su título, boletín, referencia y enlace público. & \texttt{hierarchical-retrieval} \\
RF11 & Reformular la pregunta en terminología administrativa y combinar búsqueda semántica y léxica (fusión RRF) en la pipeline jerárquica. & \texttt{hierarchical-retrieval} \\
RF12 & Responder en lenguaje claro y advertir cuando los documentos no contienen la respuesta. & \texttt{hierarchical-retrieval} \\
RF13 & Adaptar las respuestas a los temas de interés y localidades del usuario. & \texttt{contexto-personalizado} \\
RF14 & Sugerir temas y localidades a partir de las conversaciones recientes del usuario. & \texttt{procesamiento-historico} \\
\bottomrule
\end{tabularx}
\end{table}
\begin{table}[htbp]
\centering\small
\caption{Requisitos funcionales de conversación y cuenta}
\label{tab:rf-chat}
\begin{tabularx}{\textwidth}{@{}l X l@{}}
\toprule
\textbf{Id.} & \textbf{Requisito} & \textbf{Capacidad} \\
\midrule
RF15 & Guardar, listar, reabrir y borrar las conversaciones de cada usuario. & \texttt{chat-history} \\
RF16 & Mantener varias conversaciones en paralelo sin mezclar sus respuestas y recuperar la conexión antes de enviar. & \texttt{chat-ui-connection} \\
RF17 & Iniciar sesión con Auth0 dentro de la organización del ayuntamiento. & \texttt{user-auth-auth0} \\
RF18 & Gestionar el perfil del usuario y sus preferencias por organización. & \texttt{multi-tenant-user-profiles} \\
RF19 & Cerrar la sesión tras 30 minutos sin actividad, sin cortar una respuesta en curso. & \texttt{session-timeout} \\
RF20 & Ofrecer en el móvil la misma navegación que en el escritorio. & \texttt{mobile-navigation} \\
\bottomrule
\end{tabularx}
\end{table}
\begin{table}[htbp]
\centering\small
\caption{Requisitos funcionales de evaluación (Arena)}
\label{tab:rf-arena}
\begin{tabularx}{\textwidth}{@{}l X l@{}}
\toprule
\textbf{Id.} & \textbf{Requisito} & \textbf{Capacidad} \\
\midrule
RF21 & Ejecutar las dos pipelines a la vez sobre la misma pregunta y mostrar las respuestas de forma anónima (\emph{Alfa} y \emph{Beta}), en orden aleatorio. & \texttt{question-arena} \\
RF22 & Registrar el voto del usuario y revelar después qué pipeline generó cada respuesta en la página de la Arena. & \texttt{question-arena} \\
RF23 & Usar la pipeline con más victorias para responder en el chat normal. & \texttt{query-api-semantic-kernel} \\
RF24 & Calcular la tasa de victorias de cada pipeline a partir de los votos. & \texttt{ingestion-telemetry} \\
\bottomrule
\end{tabularx}
\end{table}
\begin{table}[htbp]
\centering\small
\caption{Requisitos funcionales de administración}
\label{tab:rf-admin}
\begin{tabularx}{\textwidth}{@{}l X l@{}}
\toprule
\textbf{Id.} & \textbf{Requisito} & \textbf{Capacidad} \\
\midrule
RF25 & Elegir el proveedor y el modelo de IA con el que cada ayuntamiento genera las respuestas del chat (Ollama, OpenAI, OpenRouter o Google) y guardar su clave de API. Los \emph{embeddings} los genera siempre la plataforma (sección~\ref{sec:clave-propia}). & \texttt{user-api-keys-vault} \\
RF26 & Gestionar las reglas de filtrado y lanzar descargas desde la interfaz. & \texttt{scraper-dynamic-filtering} \\
RF27 & Listar e inspeccionar los documentos almacenados, procesar uno concreto y mantener el estado de los trabajos. & \texttt{ingestion-administration} \\
RF28 & Notificar en tiempo real el progreso de la ingesta y del \emph{scraper} a los administradores. & \texttt{realtime-notifications} \\
RF29 & Mostrar el consumo de IA de la organización, el coste de ingesta de la plataforma y los resultados de la Arena, según el rol. & \texttt{admin-metrics-dashboard} \\
\bottomrule
\end{tabularx}
\end{table}

\section{Requisitos no funcionales}
\label{sec:rnf}

% TODO: añadir los requisitos de latencia con cifras cuando estén medidas (capítulo de pruebas).
Los requisitos no funcionales proceden también de las especificaciones, principalmente de las que se reescribieron tras la auditoría de código de octubre de 2026 (sección~\ref{sec:sdd}).

{\sloppy
\begin{description}
  \item[RNF1. Veracidad.] Las respuestas se basan solo en los documentos recuperados y los citan. Las instrucciones del modelo le piden que, si los documentos no responden a la pregunta, lo diga en lugar de responder con su propio conocimiento (\texttt{hierarchical-retrieval}, \texttt{query-api-semantic-kernel}).
  \item[RNF2. Aislamiento entre ayuntamientos.] Las conversaciones, votos, preferencias y configuración de IA de cada ayuntamiento quedan aisladas mediante filtros globales de consulta. Toda petición, llamada al \emph{hub} y tarea en segundo plano lleva su ayuntamiento, y fuera del entorno de desarrollo se rechazan las peticiones sin él (\texttt{multi-tenant-user-profiles}).
  \item[RNF3. Autorización en el servidor.] Cada operación administrativa exige el rol adecuado en la API, no solo en la interfaz; los servicios internos gRPC exigen una clave compartida; y los clientes nunca reciben detalles internos de los errores (\texttt{api-authorization}).
  \item[RNF4. Protección de credenciales.] Las claves de API de cada ayuntamiento se guardan cifradas, nunca se devuelven al cliente y solo se envían a su propio proveedor. Se valida siempre el certificado TLS (\texttt{user-api-keys-vault}, \texttt{llm-provider-security}).
  \item[RNF5. Tolerancia a fallos de los proveedores de IA.] Cuando el proveedor del ayuntamiento agota su cuota (HTTP 429) antes de generar texto, la pipeline \emph{baseline} responde con un modelo de respaldo de la plataforma, con sus propias credenciales, y lo avisa al usuario. El resto de errores del modelo se devuelven como mensajes tipados, sin detalles del proveedor (\texttt{query-api-semantic-kernel}).
  \item[RNF6. Fiabilidad de la ingesta.] Los mensajes de ingesta son persistentes y confirmados por el \emph{broker}; los fallos se reintentan un número limitado de veces (tres por defecto) y después pasan a una cola de mensajes fallidos. Los consumidores sobreviven a las caídas del \emph{broker} (\texttt{messaging-reliability}).
  \item[RNF7. Consistencia de datos.] PostgreSQL es la fuente de verdad y se confirma antes de escribir los vectores; la reingesta es idempotente y no deja vectores huérfanos (\texttt{ingestion-consistency}).
  \item[RNF8. Acceso respetuoso a las fuentes.] El \emph{scraper} limita la frecuencia de peticiones a cada boletín (2 por segundo al BOE y 1 por segundo al BOJA), reintenta hasta tres veces con espera exponencial ante errores de red, 429 y 5xx, y se detiene de forma ordenada (\texttt{go-scraper}).
  \item[RNF9. Observabilidad y privacidad.] Todos los servicios exportan errores, trazas, métricas y registros con OpenTelemetry, con trazas de extremo a extremo; las métricas de uso, coste y latencia de los modelos nunca incluyen el texto de las preguntas ni de las respuestas (\texttt{observability-grafana}).
  \item[RNF10. Mantenibilidad.] El código sigue una arquitectura limpia cuyas reglas de dependencia verifican pruebas automáticas, y la integración continua cubre todos los lenguajes del proyecto (\texttt{clean-architecture-boundaries}, \texttt{build-quality-gates}).
  \item[RNF11. Accesibilidad y uso en el móvil.] En el móvil, el menú lateral es un diálogo modal accesible: mueve el foco a su primer control, lo mantiene dentro mientras está abierto, expone su estado a los lectores de pantalla y se cierra con la tecla Escape, al navegar o al pulsar fuera (\texttt{mobile-navigation}).
\end{description}
\par}

\section{Casos de uso}
\label{sec:casos-uso}

Los casos de uso describen qué hace cada actor con el sistema, sin entrar en cómo está construido. La figura~\ref{fig:casos-uso} los reúne en un diagrama general y la tabla~\ref{tab:casos-uso} los relaciona con los requisitos funcionales que cubren. El administrador de organización hereda los casos de uso del usuario municipal, porque es un usuario más de su ayuntamiento con permisos adicionales. Los casos CU1 y CU11 incluyen el procesamiento de documentos de CU12, que es el mismo tanto si la descarga es diaria como si es bajo demanda.

% Bibliotecas de TikZ para las elipses y las flechas UML del diagrama.
\usetikzlibrary{shapes.geometric,arrows.meta,calc}
\begin{figure}[htbp]
\centering
\resizebox{\textwidth}{!}{%
\begin{tikzpicture}[
    font=\footnotesize,
    uc/.style={draw, ellipse, minimum width=4.6cm, minimum height=0.8cm, inner sep=0pt, align=center, execute at begin node=\setlength{\baselineskip}{1.15em}},
    lnk/.style={thin},
    dep/.style={thin, dashed, -{Stealth}},
    gen/.style={thin, -{Triangle[open]}}
  ]
  % Actor como figura de palitos: \actor{nombre}{posición}{etiqueta}
  \newcommand{\actor}[3]{%
    \begin{scope}[shift={#2}]
      \draw (0,0.45) circle (0.15);
      \draw (0,0.3) -- (0,-0.15) (-0.25,0.15) -- (0.25,0.15)
            (0,-0.15) -- (-0.2,-0.5) (0,-0.15) -- (0.2,-0.5);
      \coordinate (#1) at (0,0);
      \node[align=center, below, execute at begin node=\setlength{\baselineskip}{1.15em}] at (0,-0.5) {#3};
    \end{scope}}

  % Límite del sistema
  \draw[rounded corners] (0.2,-0.6) rectangle (12.6,-11.0);
  \node[anchor=north west, font=\footnotesize\bfseries] at (0.3,-0.65) {Asistente de boletines};

  % Usuario municipal
  \node[uc] (cu2) at (3.1,-1.6) {CU2. Consultar los boletines};
  \node[uc] (cu3) at (3.1,-2.7) {CU3. Gestionar las conversaciones};
  \node[uc] (cu4) at (3.1,-3.8) {CU4. Comparar respuestas\\en modo Arena y votar};
  \node[uc] (cu5) at (3.1,-4.9) {CU5. Mantener los temas de interés};
  % Administrador de organización
  \node[uc] (cu6) at (3.1,-6.2) {CU6. Configurar el modelo de IA};
  \node[uc] (cu7) at (3.1,-7.3) {CU7. Comparar pipelines\\en la página de la Arena};
  \node[uc] (cu8) at (3.1,-8.4) {CU8. Consultar el consumo de IA};
  % Sistema programado
  \node[uc] (cu1) at (3.1,-9.9) {CU1. Descargar los boletines\\del día anterior};
  % Administrador de plataforma
  \node[uc] (cu13) at (9.7,-4.6) {CU13. Consultar el coste\\de ingesta y la Arena};
  \node[uc] (cu9)  at (9.7,-5.8) {CU9. Gestionar las reglas\\de filtrado};
  \node[uc] (cu10) at (9.7,-7.0) {CU10. Administrar\\los documentos};
  \node[uc] (cu11) at (9.7,-8.2) {CU11. Lanzar una descarga\\bajo demanda};
  % Caso incluido
  \node[uc] (cu12) at (9.7,-9.9) {CU12. Procesar un documento};

  \actor{usr}{(-1.4,-3.3)}{Usuario\\municipal}
  \actor{org}{(-1.4,-7.3)}{Administrador\\de organización}
  \actor{sis}{(-1.4,-10.0)}{Sistema\\programado}
  \actor{plt}{(14.3,-6.4)}{Administrador\\de plataforma}

  \foreach \c in {cu2,cu3,cu4,cu5} \draw[lnk] ($(usr)+(0.3,0.15)$) -- (\c.west);
  \foreach \c in {cu6,cu7,cu8} \draw[lnk] ($(org)+(0.3,0.15)$) -- (\c.west);
  \foreach \c in {cu13,cu9,cu10,cu11} \draw[lnk] ($(plt)+(-0.3,0.15)$) -- (\c.east);
  \draw[lnk] ($(sis)+(0.3,0.15)$) -- (cu1.west);
  \draw[gen] (-1.4,-6.65) -- (-1.4,-4.6);

  \draw[dep] (cu1.east) -- node[above, font=\scriptsize] {«include»} (cu12.west);
  \draw[dep] (cu11.south) -- node[right, font=\scriptsize] {«include»} (cu12.north);
\end{tikzpicture}}
\caption{Diagrama general de casos de uso}
\label{fig:casos-uso}
\end{figure}

\begin{table}[htbp]
\centering\small
\caption{Casos de uso y requisitos funcionales que cubren}
\label{tab:casos-uso}
\begin{tabularx}{\textwidth}{@{}l X l l@{}}
\toprule
\textbf{Id.} & \textbf{Caso de uso} & \textbf{Actor principal} & \textbf{Requisitos} \\
\midrule
CU1  & Descargar los boletines del día anterior & Sistema programado & RF1, RF3, RF4 \\
CU2  & Consultar los boletines & Usuario municipal & RF8--RF13, RF16, RF23 \\
CU3  & Gestionar las conversaciones & Usuario municipal & RF15, RF16 \\
CU4  & Comparar respuestas en modo Arena y votar & Usuario municipal & RF21, RF22 \\
CU5  & Mantener los temas de interés & Usuario municipal & RF13, RF14, RF18 \\
CU6  & Configurar el modelo de IA & Admin.\ de organización & RF25 \\
CU7  & Comparar pipelines en la página de la Arena & Admin.\ de organización & RF21, RF22 \\
CU8  & Consultar el consumo de IA & Admin.\ de organización & RF29 \\
CU9  & Gestionar las reglas de filtrado & Admin.\ de plataforma & RF3, RF26 \\
CU10 & Administrar los documentos & Admin.\ de plataforma & RF7, RF27, RF28 \\
CU11 & Lanzar una descarga bajo demanda & Admin.\ de plataforma & RF2, RF26, RF28 \\
CU12 & Procesar un documento & (incluido) & RF5, RF6 \\
CU13 & Consultar el coste de ingesta y los resultados de la Arena & Admin.\ de plataforma & RF24, RF29 \\
\bottomrule
\end{tabularx}
\end{table}

El inicio y el cierre de sesión (RF17 y RF19) no aparecen como casos de uso propios: son una precondición de todos los casos de los actores humanos y los resuelve Auth0. La navegación en el móvil (RF20) tampoco es un caso de uso, sino una forma de acceder a todos ellos.

A continuación se detallan los cuatro casos de uso que mejor representan el sistema: uno por cada actor humano y el de la ingesta diaria. Cada tabla recoge las precondiciones, el flujo principal, los flujos alternativos y el estado en que queda el sistema, tomados de los escenarios de las especificaciones de OpenSpec.

% Plantilla de caso de uso detallado: una fila por campo y listas compactas para los flujos.
\newcommand{\cucampo}[2]{\textbf{#1} & #2 \\ \midrule}
\newcounter{cupaso}
\newenvironment{cupasos}
  {\begin{list}{\arabic{cupaso}.}{\usecounter{cupaso}\setlength{\leftmargin}{1.5em}\setlength{\labelwidth}{1.2em}%
    \setlength{\topsep}{0pt}\setlength{\partopsep}{0pt}\setlength{\itemsep}{0pt}\setlength{\parsep}{0pt}}}
  {\end{list}}
\newenvironment{cualternativos}
  {\begin{list}{}{\setlength{\leftmargin}{2.2em}\setlength{\labelwidth}{2em}%
    \setlength{\topsep}{0pt}\setlength{\partopsep}{0pt}\setlength{\itemsep}{0pt}\setlength{\parsep}{0pt}}}
  {\end{list}}

\subsection{CU1. Descargar los boletines del día anterior}
\label{sec:cu1}

\begin{table}[htbp]
\centering\small
\caption{Caso de uso CU1: descargar los boletines del día anterior}
\label{tab:cu1}
\begin{tabularx}{\textwidth}{@{}>{\raggedright\arraybackslash}p{2.6cm} X@{}}
\toprule
\cucampo{Actor}{Sistema programado (planificador del \emph{scraper}).}
\cucampo{Objetivo}{Incorporar a la base de conocimiento los documentos del BOE y del BOJA publicados el día anterior, sin intervención humana.}
\cucampo{Precondiciones}{Hay reglas de filtrado activas para los boletines programados y ningún otro trabajo de descarga en curso.}
\cucampo{Flujo principal}{%
\begin{cupasos}
  \item A la hora configurada, el planificador lanza un trabajo para el día anterior.
  \item El \emph{scraper} pide a la API las reglas de filtrado activas.
  \item El \emph{scraper} obtiene el sumario de cada boletín y descarta los documentos que no cumplen las reglas.
  \item Por cada documento que se conserva, guarda el original y el texto estructurado en el almacenamiento S3.
  \item Publica un mensaje de ingesta por documento en la cola.
  \item Se procesa cada documento (CU12).
\end{cupasos}}
\cucampo{Flujos alternativos}{%
\begin{cualternativos}
  \item[2a.] Las reglas no se pueden obtener: el trabajo termina como fallido y no se descarga ningún documento.
  \item[3a.] Ese día no hay boletín: el trabajo termina sin documentos y sin error.
  \item[3b.] Un sumario no se puede leer: se registra el error y se continúa con el día siguiente.
  \item[4a.] El almacenamiento falla: el documento no se encola y se contabiliza como error.
  \item[*a.] Ya hay otro trabajo en curso: el planificador omite esta ejecución.
\end{cualternativos}}
\cucampo{Postcondiciones}{Los documentos del día que cumplen las reglas están guardados en S3 y encolados para su procesamiento.}
\textbf{Requisitos} & RF1, RF3, RF4; RNF6, RNF8. \\
\bottomrule
\end{tabularx}
\end{table}

\subsection{CU2. Consultar los boletines}
\label{sec:cu2}

Es el caso de uso principal del sistema. La figura~\ref{fig:secuencia-chat} del capítulo de diseño muestra su secuencia de mensajes entre componentes.

\begin{table}[htbp]
\centering\small
\caption{Caso de uso CU2: consultar los boletines}
\label{tab:cu2}
\begin{tabularx}{\textwidth}{@{}>{\raggedright\arraybackslash}p{2.6cm} X@{}}
\toprule
\cucampo{Actor}{Usuario municipal.}
\cucampo{Objetivo}{Obtener una respuesta en lenguaje claro, basada en los boletines y con sus fuentes, a una pregunta formulada en lenguaje natural.}
\cucampo{Precondiciones}{El usuario ha iniciado sesión en la organización de su ayuntamiento y tiene una conversación abierta, nueva o recuperada (CU3).}
\cucampo{Flujo principal}{%
\begin{cupasos}
  \item El usuario escribe la pregunta y la envía.
  \item El sistema guarda la pregunta en la conversación.
  \item El sistema elige la pipeline con más victorias en la Arena.
  \item La pipeline recupera los fragmentos más relevantes de los boletines.
  \item El modelo de lenguaje genera la respuesta a partir de esos fragmentos y, en la pipeline \emph{baseline}, de los temas de interés y las localidades del usuario.
  \item El sistema muestra la respuesta a medida que se genera, con las referencias y los enlaces de los documentos citados.
  \item El sistema guarda la respuesta completa en la conversación.
\end{cupasos}}
\cucampo{Flujos alternativos}{%
\begin{cualternativos}
  \item[1a.] La conexión con el servidor se ha perdido: el sistema la recupera antes de enviar; si no lo consigue, avisa al usuario y no envía la pregunta.
  \item[4a.] No se recupera ningún fragmento: la pipeline jerárquica responde con un mensaje fijo que indica que no ha encontrado normativa relevante; la \emph{baseline} responde como en una conversación general, sin citar documentos, lo que cubre por ejemplo los saludos.
  \item[5a.] El proveedor de IA del ayuntamiento ha agotado su cuota antes de generar texto: la pipeline \emph{baseline} responde con el modelo de respaldo de la plataforma y lo indica.
  \item[5b.] El modelo falla: el usuario recibe un mensaje de error comprensible, sin detalles del proveedor; si ya se había mostrado parte de la respuesta, se conserva con el aviso.
  \item[6a.] El usuario cancela la respuesta: se guarda el texto recibido hasta ese momento.
\end{cualternativos}}
\cucampo{Postcondiciones}{La pregunta y la respuesta quedan guardadas en la conversación del usuario, dentro de su ayuntamiento, y la llamada al modelo queda registrada para las métricas de consumo.}
\textbf{Requisitos} & RF8--RF13, RF16, RF23; RNF1, RNF2, RNF5. \\
\bottomrule
\end{tabularx}
\end{table}

\subsection{CU6. Configurar el modelo de IA}
\label{sec:cu6}

\begin{table}[htbp]
\centering\small
\caption{Caso de uso CU6: configurar el modelo de IA}
\label{tab:cu6}
\begin{tabularx}{\textwidth}{@{}>{\raggedright\arraybackslash}p{2.6cm} X@{}}
\toprule
\cucampo{Actor}{Administrador de organización.}
\cucampo{Objetivo}{Elegir el proveedor y el modelo con el que se generan las respuestas del chat de su ayuntamiento y, si hace falta, la clave de API con la que se pagan.}
\cucampo{Precondiciones}{El usuario ha iniciado sesión con el rol \texttt{admin-organizacion}.}
\cucampo{Flujo principal}{%
\begin{cupasos}
  \item El administrador abre la página de configuración.
  \item El sistema muestra el proveedor, el modelo, la temperatura y el punto de acceso actuales, e indica si hay una clave guardada sin mostrarla.
  \item El administrador elige el proveedor (Ollama, OpenAI, OpenRouter o Google) y el modelo, y escribe la clave de API.
  \item El administrador guarda los cambios.
  \item El sistema cifra la clave y guarda la configuración de su ayuntamiento.
\end{cupasos}}
\cucampo{Flujos alternativos}{%
\begin{cualternativos}
  \item[3a.] El administrador no escribe una clave nueva: se conserva la que estaba guardada.
  \item[3b.] Elige un proveedor que necesita clave y no hay ninguna: las respuestas del chat fallan con un mensaje que indica que el modelo no está configurado.
  \item[*a.] El usuario no tiene el rol de administrador de organización: la página no está disponible y la API rechaza la petición.
\end{cualternativos}}
\cucampo{Postcondiciones}{Las respuestas de la pipeline \emph{baseline} de ese ayuntamiento se generan con el modelo elegido. La clave nunca se devuelve al cliente y solo se envía a su proveedor. Los \emph{embeddings} y la pipeline jerárquica siguen usando los modelos de la plataforma (sección~\ref{sec:clave-propia}).}
\textbf{Requisitos} & RF25; RNF2, RNF3, RNF4. \\
\bottomrule
\end{tabularx}
\end{table}

\subsection{CU11. Lanzar una descarga bajo demanda}
\label{sec:cu11}

\begin{table}[htbp]
\centering\small
\caption{Caso de uso CU11: lanzar una descarga bajo demanda}
\label{tab:cu11}
\begin{tabularx}{\textwidth}{@{}>{\raggedright\arraybackslash}p{2.6cm} X@{}}
\toprule
\cucampo{Actor}{Administrador de plataforma.}
\cucampo{Objetivo}{Incorporar a la base de conocimiento los documentos de un boletín para un rango de fechas, por ejemplo para cargar días anteriores a la puesta en marcha o recuperar un día que falló.}
\cucampo{Precondiciones}{El usuario ha iniciado sesión con el rol \texttt{admin-plataforma}.}
\cucampo{Flujo principal}{%
\begin{cupasos}
  \item El administrador abre la página de reglas de filtrado.
  \item Elige el boletín, el rango de fechas y, opcionalmente, las secciones.
  \item El administrador lanza la descarga.
  \item El sistema pide al \emph{scraper} que descargue el rango y sigue los pasos 2 a 5 de CU1 para cada día.
  \item El sistema notifica en tiempo real al administrador el resultado y el número de documentos guardados.
  \item Se procesa cada documento (CU12).
\end{cupasos}}
\cucampo{Flujos alternativos}{%
\begin{cualternativos}
  \item[2a.] El rango no es válido (el final es anterior al inicio o posterior a hoy, o supera los 31 días): el sistema rechaza la descarga e indica el motivo.
  \item[3a.] Ya hay una descarga en curso: el sistema la rechaza sin llamar al \emph{scraper}.
  \item[4a.] El \emph{scraper} no puede obtener las reglas de filtrado: la descarga termina como fallida, sin descargar ningún documento, y se notifica al administrador.
\end{cualternativos}}
\cucampo{Postcondiciones}{Los documentos del rango que cumplen las reglas están guardados en S3 y encolados para su procesamiento, y su progreso puede seguirse desde la página de documentos (CU10).}
\textbf{Requisitos} & RF2, RF3, RF26, RF28; RNF3, RNF8. \\
\bottomrule
\end{tabularx}
\end{table}

% ======================================================================
% Capítulo: Diseño del sistema
% Fase 2 del anteproyecto. Extensión orientativa: 8--10 páginas.
% ======================================================================
\chapter{Diseño del sistema}
\label{cap:diseno}

\section{Arquitectura general}
\label{sec:arquitectura}
% TODO: Los tres componentes desacoplados (Go, .NET, Angular), gateway, workers y almacenes.

\section{Arquitectura limpia en .NET}
\label{sec:clean}
% TODO: Capas, puertos y adaptadores, tests de arquitectura.

\section{Diagramas UML}
\label{sec:uml}
% TODO: Clases de los módulos clave y secuencias de pregunta en el chat e ingesta de un boletín.

\section{Modelo de datos híbrido}
\label{sec:modelo-datos}
% TODO: Diagrama E-R de PostgreSQL y colecciones de Qdrant (\texttt{document\_chunks}, \texttt{child\_fragments}).

\section{Indexación vectorial}
\label{sec:indexacion}
% TODO: Modelos de embeddings y dimensiones, similitud coseno, HNSW, límite de pgvector y migración a Qdrant (ADR 0001).

\section{Flujo de ingesta asíncrona}
\label{sec:ingesta-diseno}
% TODO: Contrato de RabbitMQ, colas de reintento y DLQ, idempotencia y estados del documento.

\section{Pipelines RAG}
\label{sec:pipelines}
% TODO: Baseline (chunks fijos, búsqueda densa) frente a jerárquica (expansión de consulta, híbrida RRF, contexto padre, lenguaje claro).

\section{Abstracción del proveedor de IA y control de coste}
\label{sec:proveedor-ia}
% TODO: Semantic Kernel, fábrica de modelos por tenant, fallback ante 429 y registro de coste por llamada.

\section{Tiempo real y seguridad}
\label{sec:seguridad}
% TODO: SignalR, multi-tenancy con Auth0 (acceso por ayuntamiento), roles y cifrado de claves.

% ======================================================================
% Capítulo: Implementación
% Fases 3, 5, 6 y 7 del anteproyecto. Extensión orientativa: 8--10 páginas. Código corto solo donde aporte; el resto, a apéndices.
% ======================================================================
\chapter{Implementación}
\label{cap:implementacion}

\section{Módulo de ingesta en Go}
\label{sec:impl-go}
% TODO: Clientes de BOE y BOJA, descarga concurrente, extracción, S3, eventos y filtros por gRPC.

\section{Backend .NET: ingesta}
\label{sec:impl-ingesta}
% TODO: Consumidores, fragmentación de BOE y BOJA, municipio y fecha, embeddings, Postgres y Qdrant.

\section{Orquestación RAG}
\label{sec:impl-rag}
% TODO: Recuperación híbrida, prompts (\texttt{PromptCatalog}), citas de fuentes y streaming por SignalR.

\section{Contexto personalizado y Arena}
\label{sec:impl-arena}
% TODO: Preferencias, comparación de pipelines, votos, \emph{shadow testing}; la Arena como herramienta de validación con funcionarios.

\section{Frontend Angular}
\label{sec:impl-angular}
% TODO: Organización por features, señales, chat, visor de citas y adaptación a móvil.

\section{Orquestación y despliegue}
\label{sec:impl-despliegue}
% TODO: AppHost de Aspire, Docker Compose, Dokploy y CI/CD.

\section{Observabilidad y coste}
\label{sec:impl-observabilidad}
% TODO: OpenTelemetry, métricas GenAI (tokens, latencia, TTFT, coste en EUR), Grafana y privacidad (ADR 0002). Tabla de coste por pregunta: modelos frontera frente a locales.

\section{Auditoría y refactorización}
\label{sec:impl-auditoria}
% TODO: Auditoría de código de octubre de 2026 y sus fases de corrección.

% ======================================================================
% Capítulo: Estudio comparativo de estrategias de \emph{chunking}
% Fase 4 del anteproyecto. Extensión orientativa: 5--7 páginas. Resultados completos en el apéndice~\ref{ap:chunking}.
% ======================================================================
\chapter{Estudio comparativo de estrategias de \emph{chunking}}
\label{cap:estudio-chunking}

\section{Estrategias comparadas}
\label{sec:estrategias}
% TODO: Fixed-size con solapamiento, recursive, semantic y jerárquica padre/hijo, con sus parámetros.

\section{Diseño experimental}
\label{sec:diseno-experimental}
% TODO: Corpus fijo, preguntas de referencia con \emph{ground truth}, mismos modelos de embeddings y generación.

\section{Entorno de evaluación local}
\label{sec:entorno-local}
% TODO: RAGAS en local con Ollama como juez; justificación por coste.

\section{Métricas}
\label{sec:metricas}
% TODO: Precisión y recall de contexto, fidelidad, relevancia, tasa de alucinación, latencia y coste de indexación.

\section{Resultados}
\label{sec:resultados}
% TODO: Tablas y gráficas por estrategia y métrica.

\section{Discusión}
\label{sec:discusion}
% TODO: Qué estrategia gana, con qué coste, y si confirma la jerárquica en producción.

\section{Amenazas a la validez}
\label{sec:validez}
% TODO: Tamaño del conjunto de preguntas, modelo local como juez, representatividad del corpus.

% ======================================================================
% Capítulo: Pruebas y validación
% Fases 7 y 8 del anteproyecto. Extensión orientativa: 4--6 páginas.
% ======================================================================
\chapter{Pruebas y validación}
\label{cap:pruebas}

\section{Estrategia de pruebas}
\label{sec:estrategia-pruebas}
% TODO: Unitarias, integración con Testcontainers, arquitectura, Go con \texttt{-race} y Angular.

\section{Pruebas de seguridad y aislamiento}
\label{sec:pruebas-seguridad}
% TODO: Tenants, autorización y privacidad de la telemetría.

\section{Validación con usuarios reales}
\label{sec:validacion-usuarios}
% TODO: Protocolo con funcionarios vía Arena (cada ayuntamiento como tenant). Resultados si llega el acceso; si no, validación preparada y escenarios simulados.

\section{Rendimiento}
\label{sec:rendimiento}
% TODO: Latencia del chat, tiempo de ingesta por documento y tasa de errores.

% ======================================================================
% Capítulo: Conclusiones y líneas futuras
% Extensión orientativa: 3--4 páginas.
% ======================================================================
\chapter{Conclusiones y líneas futuras}
\label{cap:conclusiones}

\section{Cumplimiento de objetivos}
\label{sec:cumplimiento}
% TODO: OE1--OE7 uno a uno, señalando cambios de alcance.

\section{Aportaciones}
\label{sec:aportaciones}
% TODO: Arena y multi-tenancy como base de validación; métricas de coste para elegir modelo.

\section{Lecciones aprendidas}
\label{sec:lecciones}
% TODO: Técnicas y de proceso (SDD con agentes de IA).

\section{Líneas futuras}
\label{sec:futuras}
% TODO: Boletines provinciales, alertas de subvenciones y plazos, observabilidad de nivel 2, apertura a la ciudadanía.

\paragraph{Clave de API propia en producción.} La plataforma permite que cada ayuntamiento configure su proveedor y su clave de API para el chat, mientras la ingesta y los \emph{embeddings} se mantienen centralizados (sección~\ref{sec:clave-propia}). Hoy no se usa por los costes; activarlo permitiría que cada entidad pague solo sus conversaciones. Quedan como mejoras posibles llevar las claves a un almacén de secretos externo y valorar si compensa pasar los \emph{embeddings} de la plataforma a un proveedor externo, lo que exigiría reindexar todos los boletines con el nuevo modelo.


%────────────────────── Apéndices ──────────────────────────────────────
\begin{umaappendices}
\umaappendix{Resultados completos del estudio de \emph{chunking}}
\label{ap:chunking}

Este apéndice recoge los resultados por pregunta de la ejecución \texttt{20261023-112730} (cuatro variantes, \emph{prompts} ajustados, RTX 4060) que resume el capítulo~\ref{cap:estudio-chunking}. Los ficheros completos de las ejecuciones (respuestas, contextos, puntuaciones y latencias) están en el repositorio, en \texttt{eval/runs}.

\section{Fidelidad y recall de contexto por pregunta}

La tabla~\ref{tab:ap-por-pregunta} da, para cada pregunta, la fidelidad (F) y el recall de contexto (CR) de cada variante. Un guion indica que el juez no pudo puntuar la muestra. En las tres preguntas sin respuesta en el corpus (q23 a q25) el recall de contexto mide la cobertura de una referencia que dice que los boletines no recogen la información.

\begin{small}
\begin{longtable}{@{}l l rr rr rr rr@{}}
\caption{Fidelidad (F) y recall de contexto (CR) por pregunta y variante}\label{tab:ap-por-pregunta}\\
\toprule
& & \multicolumn{2}{c}{Tamaño fijo} & \multicolumn{2}{c}{Recursiva} & \multicolumn{2}{c}{Semántica} & \multicolumn{2}{c}{Jerárquica} \\
\cmidrule(lr){3-4}\cmidrule(lr){5-6}\cmidrule(lr){7-8}\cmidrule(lr){9-10}
\textbf{Id.} & \textbf{Categoría} & F & CR & F & CR & F & CR & F & CR \\
\midrule
\endfirsthead
\toprule
\textbf{Id.} & \textbf{Categoría} & F & CR & F & CR & F & CR & F & CR \\
\midrule
\endhead

q01 & empleo-publico & 1,00 & 1,00 & 0,80 & 1,00 & 1,00 & 1,00 & 0,25 & 1,00 \\
q02 & empleo-publico & 1,00 & 1,00 & 0,00 & 0,67 & 1,00 & 1,00 & 1,00 & 1,00 \\
q03 & empleo-publico & 0,43 & 1,00 & 0,75 & 1,00 & 1,00 & 0,67 & 1,00 & 1,00 \\
q04 & empleo-publico & 0,75 & 1,00 & 0,75 & 1,00 & 1,00 & 1,00 & 0,83 & 1,00 \\
q05 & empleo-publico & 0,50 & 1,00 & 0,00 & 1,00 & 1,00 & 1,00 & 0,50 & 1,00 \\
q06 & empleo-publico & 0,33 & 0,00 & 1,00 & 1,00 & 1,00 & 0,00 & 0,83 & 1,00 \\
q07 & empleo-publico & 1,00 & 1,00 & 0,67 & 1,00 & 0,67 & 1,00 & 0,78 & 1,00 \\
q08 & empleo-publico & 1,00 & 1,00 & 1,00 & 1,00 & 0,50 & 1,00 & 0,88 & 1,00 \\
q09 & empleo-publico & 1,00 & 1,00 & 1,00 & 0,00 & 0,50 & 1,00 & 0,50 & 1,00 \\
q10 & empleo-publico & 1,00 & 1,00 & 1,00 & 1,00 & 1,00 & 1,00 & 0,80 & 1,00 \\
q11 & ayudas & 0,75 & 0,50 & 1,00 & 0,50 & 0,86 & 0,50 & 1,00 & 0,50 \\
q12 & ayudas & 1,00 & 0,00 & 1,00 & 0,67 & 1,00 & 1,00 & 1,00 & 1,00 \\
q13 & ayudas & 1,00 & 1,00 & 1,00 & 1,00 & 0,67 & 0,00 & 1,00 & 1,00 \\
q14 & ayudas & 0,71 & 0,00 & 1,00 & 1,00 & 1,00 & 1,00 & --- & 1,00 \\
q15 & ayudas & 1,00 & 0,00 & 0,75 & 0,50 & 1,00 & 0,00 & 1,00 & 0,00 \\
q16 & ayudas & 0,00 & 1,00 & 1,00 & 1,00 & 0,00 & 1,00 & 1,00 & 1,00 \\
q17 & normativa & 0,67 & 1,00 & 0,67 & 1,00 & 0,67 & 0,00 & 1,00 & 1,00 \\
q18 & normativa & 1,00 & 1,00 & 1,00 & 1,00 & 0,50 & 1,00 & 0,67 & 1,00 \\
q19 & normativa & 0,25 & 0,00 & 1,00 & 1,00 & 1,00 & 1,00 & 0,80 & 1,00 \\
q20 & normativa & 1,00 & 1,00 & 1,00 & 1,00 & 1,00 & 1,00 & 0,67 & 1,00 \\
q21 & tributos & 1,00 & 1,00 & 0,50 & 1,00 & 0,67 & 1,00 & 0,67 & 1,00 \\
q22 & tributos & 1,00 & 0,00 & 1,00 & 0,00 & 1,00 & 1,00 & 0,67 & 1,00 \\
q23 & sin-respuesta & 1,00 & 0,00 & 1,00 & 1,00 & 0,67 & 1,00 & 0,50 & 1,00 \\
q24 & sin-respuesta & 0,33 & 0,00 & 0,71 & 0,00 & 1,00 & 0,00 & 0,43 & 1,00 \\
q25 & sin-respuesta & 0,67 & 0,00 & 1,00 & 0,50 & 1,00 & 0,50 & 0,40 & 0,50 \\
\bottomrule
\end{longtable}
\end{small}

\section{Comparación con los \emph{prompts} anteriores}

La tabla~\ref{tab:ap-prompts} compara la ejecución con los \emph{prompts} anteriores (\texttt{20261023-110829}) y la ejecución con los \emph{prompts} ajustados (\texttt{20261023-112730}), con las mismas preguntas, el mismo corpus y el mismo hardware. La recuperación es idéntica en las dos, así que las diferencias en precisión y recall de contexto se deben solo al juez y dan una idea de su ruido.

\begin{table}[htbp]
\centering\small
\caption{Medias por variante antes y después del ajuste de los \emph{prompts}}
\label{tab:ap-prompts}
\begin{tabular}{@{}l rr rr rr rr@{}}
\toprule
& \multicolumn{2}{c}{Precisión ctx.} & \multicolumn{2}{c}{Recall ctx.} & \multicolumn{2}{c}{Fidelidad} & \multicolumn{2}{c}{Alucinación} \\
\cmidrule(lr){2-3}\cmidrule(lr){4-5}\cmidrule(lr){6-7}\cmidrule(lr){8-9}
\textbf{Variante} & antes & después & antes & después & antes & después & antes & después \\
\midrule

Tamaño fijo & 0,44 & 0,44 & 0,62 & 0,62 & 0,74 & 0,78 & 56\,\% & 44\,\% \\
Recursiva & 0,53 & 0,52 & 0,81 & 0,79 & 0,85 & 0,82 & 28\,\% & 36\,\% \\
Semántica & 0,46 & 0,42 & 0,79 & 0,75 & 0,82 & 0,83 & 36\,\% & 36\,\% \\
Jerárquica & 0,77 & 0,77 & 0,88 & 0,92 & 0,68 & 0,76 & 48\,\% & 46\,\% \\
\bottomrule
\end{tabular}
\end{table}


\umaappendix{Manual de usuario}
\label{ap:manual-usuario}
% TODO: Para el personal municipal, con capturas: chat, citas, historial, Arena, perfil.

\umaappendix{Guía de despliegue}
\label{ap:despliegue}
% TODO: Para el equipo de TI, a partir de docs/setup.md y docs/deploy-dokploy.md.

\umaappendix{Contrato de mensajería y especificaciones}
\label{ap:especificaciones}
% TODO: Colas, mensajes y estados (docs/messaging-contract.md) y listado de capacidades de OpenSpec con un ejemplo completo.

\umaappendix{Prompts del sistema}
\label{ap:prompts}
% TODO: Texto literal de cada prompt de PromptCatalog.cs junto al requisito de OpenSpec que lo define.

\end{umaappendices}

%────────────────────── Back cover and bibliography ────────────────────
\backmatter
\nocite{*}  % TODO: quitar cuando los capítulos ya citen las referencias
\printbibliography[title={Bibliografía}]
\cleardoublepage
\MakeUMABackCover

\end{document}
% ======================================================================
